Soft-biometric attributes are usually judged by how accurately they can be predicted, yet identification depends on whether a predicted attribute improves the rank of the true writer. We measure this downstream utility for handwriting handedness with a paired protocol that fixes the gallery and identity scores and varies only the auxiliary information: known hand, a known random partition, hard filtering, and calibrated fixed or nested learned-weight fusion. Handedness models exclude the queried writer from training and calibration, and identity normalization uses gallery references only. On BanglaHand (42 writers, six representations, 120 settings), known handedness raises top-1 accuracy by 6.1 percentage points, whereas an arbitrary known 21/21 partition raises it by 6.9. Identity errors cluster within hand: when one impostor outranks the true writer, it has the opposite hand in 34--37\% of cases for neural matchers, against 51\% for a random impostor. With query handedness predicted at 65.5\% accuracy, hard filtering lowers top-1 accuracy by 14.3 points and nested fusion changes it by $-0.08$ points; the 90\% interval ($-0.96$ to $+0.67$) excludes recovering more than 11\% of the oracle gain. On 96 CSAFE writers, with gallery and queries from separate capture sessions, no primary fusion contrast shows a gain. A handwriting-specific SIFT-VLAD matcher (91--96\% top-1) leaves no oracle opportunity at all, and gender as a second attribute repeats the pattern: known gender is no better than a random partition, errors cluster within gender, and none of 16 pre-specified fusion contrasts shows a gain. In these cohorts, neither attribute accuracy nor a correct attribute established identification utility. The protocol separates opportunity, specificity, achieved utility, mechanism and writer-associated information, and can be applied to other auxiliary attributes.
